% NEW ACRONYM — add to formatting/Acronyms.tex:
% \DeclareAcronym{CPU}{short=CPU, long=Central Processing Unit, tag=abbrev, sort=CPU}
% \DeclareAcronym{PGS}{short=PGS, long=Projected Gauss--Seidel, tag=abbrev, sort=PGS}
% \DeclareAcronym{TGS}{short=TGS, long=Temporal Gauss--Seidel, tag=abbrev, sort=TGS}
% \DeclareAcronym{SDF}{short=SDF, long=Signed Distance Field, tag=abbrev, sort=SDF}

%===================================================================%
\section{Physics Simulation in Isaac Sim 5.1}
\label{sec:physics_sim}
%===================================================================%

Physics-based simulation in Isaac~Sim~5.1 is implemented as a layered stack that separates
(i) scene and asset description, (ii) binding of these descriptions to a concrete backend, and
(iii) numerical solvers that evolve the physical state over time.
In the Omniverse ecosystem, physics properties are authored in \ac{USD} through the
\texttt{UsdPhysics} schema (engine-agnostic physics description) and \texttt{PhysXSchema}
(backend-specific extensions), allowing assets and environments to be defined declaratively in
the scene graph~\cite{openusd_usdphysics_schema,omni_physx_overview}.
During simulation, the \texttt{omni.physx} runtime parses the \ac{USD} stage, constructs the
corresponding physics scene, advances the simulation, and writes back poses and state to
\ac{USD}~\cite{omni_physx_overview}.
Isaac~Sim~5.1 uses NVIDIA PhysX~5 as its primary physics backend, providing rigid-body
dynamics, articulations, contacts, and particle-based phenomena, with optional
\ac{GPU}-accelerated pipelines for throughput-critical workloads such as robot learning
experiments with many parallel environments~\cite{isaac_sim_simulation_fundamentals_510,physx_welcome_513,physx_gpu_rigid_bodies_510}.

A detailed understanding of the underlying physics models and numerical methods is
required for two reasons.
First, sim-to-real transfer is fundamentally limited by discrepancies between simulated and
real dynamics (the \emph{reality gap}), where unmodelled effects, incorrect parameters (e.g.,
friction and damping), and solver artefacts can lead to policies that exploit simulator-specific
regularities~\cite{Jakobi1995RealityGap,Salvato2021RealityGapSurvey}.
Second, in \ac{RL}, the simulator acts as a data-generating process; thus, modelling and solver
choices directly shape the transition distribution used for learning, influencing both policy
quality and generalization~\cite{Salvato2021RealityGapSurvey,Peng2018DynamicsRandomization}.
In practice, this implies that solver selection, time-step choice, contact modelling, and
collision geometry approximations must be treated as part of the learning system design, not
as interchangeable implementation details~\cite{Peng2018DynamicsRandomization,physx_best_practices_541}.

\begin{figure}[ht!]
\centering
\includegraphics[width=0.7\textwidth]{isaac_sim_physics_stack}
\caption[Isaac Sim physics stack]{Conceptual physics simulation stack in Isaac~Sim~5.1: physics is authored in \ac{USD} via \texttt{UsdPhysics} and \texttt{PhysXSchema}, interpreted by \texttt{omni.physx}, and simulated using PhysX~5 with optional \ac{GPU}-accelerated pipelines.}\label{fig:isaac_sim_physics_stack}
\end{figure}

\subsection{Rigid Body Dynamics}

Rigid-body simulation assumes that each object moves as a rigid transform, i.e., deformation
is neglected and the state is represented by position/orientation and linear/angular velocity.
PhysX represents the simulated world as a set of actors (static, kinematic, dynamic) with
attached collision shapes and material properties; dynamics are advanced under external
loads (e.g., gravity), contact constraints, and joint constraints~\cite{physx_welcome_513,physx_rigid_body_dynamics_511}.
In Isaac~Sim, these concepts are exposed through \ac{USD} physics schemas and their PhysX
extensions, enabling rigid bodies, colliders, and constraints to be authored on scene
primitives~\cite{isaac_sim_simulation_fundamentals_510,openusd_usdphysics_schema}.

\paragraph{Articulated bodies, joints, and robot mechanisms.}
Robotic manipulators are modelled as articulated mechanisms, i.e., rigid links connected by
joints that constrain relative motion and optionally apply actuation (drives).
PhysX provides an \emph{articulation} abstraction that is simulated in reduced coordinates,
where the configuration is parameterized by the root pose and joint coordinates rather than
the world pose of each link~\cite{physx_articulations_513}.
This formulation is designed to improve numerical fidelity for mechanisms (e.g., reduced
joint error and improved behavior under large mass ratios) compared to independently
simulating rigid bodies connected by generic constraints~\cite{physx_articulations_513}.
In robotics applications, robot mechanisms such as tendon-driven arms and multi-finger
grippers are simulated as chains of rigid links with joint constraints and joint drives,
such that control inputs act on joint coordinates while the solver enforces inter-link
constraints and contacts~\cite{physx_articulations_513,isaac_sim_simulation_fundamentals_510}.
PhysX also exposes explicit \emph{articulation tendons} as additional constraint structures
within an articulation, including fixed tendons (coupling joint coordinates) and spatial
tendons (spring-damper and limit constraints defined over attachment trees across links),
which are relevant when cable-like actuation or joint coupling must be represented at the
physics level~\cite{physx_articulations_513,omni_physics_tendon_authoring}.

\paragraph{Contacts and friction models.}
Rigid-body collision response is formulated as a constrained dynamics problem in which
contacts prevent interpenetration and can transfer or dissipate energy through restitution and
friction~\cite{physx_rigid_body_collision_510,isaac_sim_simulation_fundamentals_510}.
In PhysX, friction and restitution are specified via materials attached to shapes; friction is
parameterized by static and dynamic coefficients and can be combined for contacting pairs
using configurable combine modes~\cite{physx_rigid_body_dynamics_511,physx_material_api_510}.
Because contact-rich manipulation is common in robotic assembly, grasping, and textile
handling, parameter choices for friction and restitution are directly coupled to learning
outcomes and to sim-to-real robustness~\cite{Salvato2021RealityGapSurvey,Peng2018DynamicsRandomization}.

\paragraph{\acs{TGS} constraint solving in PhysX.}
Rigid-body dynamics with contacts and joints is solved in PhysX using iterative constraint
solvers. PhysX supports a classic \acs{PGS}-style solver and the newer \acs{TGS} solver, with
\acs{TGS} introduced in PhysX~5.1 and recommended for improved convergence and joint
drive accuracy in many scenarios~\cite{physx_rigid_body_dynamics_511}.
Both solvers operate iteratively over constraint manifolds (contacts, joints, friction), and
accuracy is traded for performance by limiting iteration counts.
A key motivation for temporal formulations is the observation that, for stiff constraint
systems, replacing multiple solver iterations at a long time step by multiple smaller substeps
can yield superior stability and accuracy at comparable computational cost~\cite{Macklin2019SmallSteps,physx_best_practices_541}.
This perspective is consistent with PhysX guidance to consider smaller time steps (or
substepping) when joint stability and contact behavior are insufficient under a fixed number
of iterations, rather than indefinitely increasing iteration counts~\cite{physx_best_practices_541}.

\paragraph{Simulation parameters and numerical tradeoffs.}
For robotics and \ac{RL}, the most consequential rigid-body parameters are those controlling
temporal discretization, solver convergence, and contact generation.
The simulation time step $h$ determines how frequently dynamics are integrated and how
often the policy interacts with the environment; overly large $h$ can lead to missed
collisions, unstable joint behavior, and biased contact impulses, whereas smaller $h$
improves accuracy at increased compute cost~\cite{physx_simulation_541,physx_best_practices_541}.
Solver iteration counts control how many passes are made over constraints; PhysX allows
per-actor iteration counts (position and velocity iterations), and the effective iterations for a
connected island are determined by the maximum requested within that island~\cite{physx_simulation_541,physx_rigid_body_dynamics_511}.
Contact generation is influenced by the pair \texttt{contactOffset} and \texttt{restOffset}: contact
points are generated once shapes are within the combined contact distance, while the rest
offset defines the target separation distance at rest, enabling predictive enforcement and
reducing jitter at the cost of introducing a small separation bias~\cite{physx_advanced_collision_detection_512}.
Finally, friction coefficients (static/dynamic) and restitution parameters are specified in
materials and substantially affect sliding, sticking, and impact behavior~\cite{physx_material_api_510,physx_rigid_body_dynamics_511}.

\begin{table}[htb]
    \centering
    \caption{Selected rigid-body simulation parameters relevant to robot learning in Isaac~Sim / PhysX.}\label{tab:physx_rigid_body_parameters}
    \begin{tabular}{l c c}
    \hline
    \textbf{Parameter} & \textbf{Location} & \textbf{Primary effect} \\
    \hline
    Time step $h$ & Physics scene & Accuracy vs.\ throughput; contact/joint stability~\cite{physx_simulation_541} \\
    Solver iterations & Actor / island & Constraint convergence; joint error; contact quality~\cite{physx_rigid_body_dynamics_511} \\
    \texttt{contactOffset}/\texttt{restOffset} & Shape & Contact prediction and resting separation~\cite{physx_advanced_collision_detection_512} \\
    Friction coefficients & Material & Sliding/sticking behavior; grasp stability~\cite{physx_material_api_510} \\
    \hline
    \end{tabular}
\end{table}

\paragraph{\ac{GPU}-accelerated rigid bodies for parallel \ac{RL} environments.}
PhysX provides a \ac{GPU} rigid-body pipeline that accelerates broad-phase, narrow-phase,
contact generation, and constraint solving on CUDA-capable hardware, while exposing the
same API concepts as the \ac{CPU} pipeline~\cite{physx_gpu_rigid_bodies_510}.
Large performance gains are expected when many rigid bodies are active, which aligns with
robot learning workloads where hundreds to thousands of environment replicas are simulated
in parallel to increase sample throughput~\cite{physx_gpu_rigid_bodies_510,Makoviychuk2021IsaacGym}.
In IsaacLab, PhysX \ac{GPU} simulation is enabled by default and can be configured
explicitly; however, \ac{GPU} simulation requires preallocation of internal buffers and may
fail when insufficient capacities are specified, which makes sizing and profiling part of
simulation configuration for large-scale training~\cite{isaaclab_sim_doc}.
The broader motivation for \ac{GPU}-resident simulation in \ac{RL} is that physics stepping
and policy learning can share device memory, avoiding \ac{CPU}--\ac{GPU} transfer
bottlenecks; this design has been demonstrated to yield substantial throughput improvements
in GPU-based robot learning platforms~\cite{Makoviychuk2021IsaacGym}.

\paragraph{Collision detection and geometry approximations.}
Collision detection cost is highly sensitive to the chosen collider representation.
Primitive shapes (e.g., spheres, boxes, capsules) provide the strongest performance and
numerical robustness because intersection tests and contact generation are simple and
well-conditioned~\cite{physx_rigid_body_collision_510}.
Convex meshes provide a practical middle ground for complex objects by enabling efficient
collision queries while approximating concave geometry; PhysX requires \emph{convex mesh
cooking} to transform input geometry into an optimized representation for collision
detection~\cite{physx_geometry_541}.
In Isaac~Sim, mesh colliders are often approximated by convex hulls by default for
performance, with more expensive approximations available when needed~\cite{isaac_sim_simulation_fundamentals_510}.
When concavity is important for dynamic objects, convex decomposition (a set of convex
parts) is typically preferred over a single triangle mesh collider because triangle-mesh
contact generation is expensive and, depending on pipeline and version, may require
additional structures.
In PhysX~5.3+, dynamic triangle meshes in the \ac{GPU} collision pipeline require an
associated \ac{SDF} to accelerate collision detection, increasing memory use but improving
query performance~\cite{physx_geometry_530}.
Therefore, collider selection becomes a modelling choice that trades geometric fidelity
against computational cost and solver stability, and this tradeoff must be balanced against
the throughput requirements of \ac{RL} training.

% NOTE: Cloth and deformable object simulation content has been moved to
% the Master Thesis background (Chapter 2, Section 2.2: Cloth and Deformable
% Object Simulation) because cloth dynamics are not used in the project thesis
% experiments, which focus exclusively on rigid-object manipulation.
% See: master_thesis/thesis/chapters/background/2_2_ClothSimulation.tex

% ---------------------------------------------------------------------------
% BibTeX keys used in this section:
%  - omni_physx_overview                    (Omni PhysX overview: USD $\leftrightarrow$ PhysX binding)
%  - openusd_usdphysics_schema              (OpenUSD UsdPhysics schema overview)
%  - isaac_sim_simulation_fundamentals_510  (Isaac Sim 5.1 physics fundamentals)
%  - physx_welcome_513                      (PhysX 5.1.3 introduction)
%  - physx_rigid_body_dynamics_511          (PhysX rigid body dynamics, TGS vs PGS, iterations)
%  - physx_rigid_body_collision_510         (PhysX rigid body collision overview)
%  - physx_material_api_510                 (PhysX material friction/restitution API)
%  - physx_simulation_541                   (PhysX simulation loop and iteration/island behavior)
%  - physx_best_practices_541               (PhysX best practices: time step vs iterations)
%  - physx_advanced_collision_detection_512 (PhysX contactOffset/restOffset semantics)
%  - physx_gpu_rigid_bodies_510             (PhysX GPU rigid bodies)
%  - physx_geometry_541                     (PhysX geometry/collision support overview)
%  - physx_geometry_530                     (PhysX 5.3 geometry: dynamic triangle meshes and SDFs on GPU)
%  - physx_articulations_513                (PhysX articulations and tendons)
%  - omni_physics_tendon_authoring          (Omni Physics tendon authoring guide)
%  - isaaclab_sim_doc                       (IsaacLab simulation configuration: GPU sim, buffers)
%  - physx_particle_system_513              (PhysX particle system: PBD solver, cloth support)
%  - physx_pbd_particlesystem_api           (PhysX PxPBDParticleSystem API reference)
%  - omni_physics_particles                 (Omni Physics particles: GPU PBD particles, cloth, deprecation)
%  - isaac_sim_physics_resources_510        (Isaac Sim physics limitations: particle cloth deprecation)
%  - Makoviychuk2021IsaacGym                (NeurIPS 2021 Isaac Gym: GPU simulation for RL)
%  - Jakobi1995RealityGap                   (Noise and the reality gap in simulation)
%  - Salvato2021RealityGapSurvey            (IEEE Access survey on sim-to-real in RL)
%  - Peng2018DynamicsRandomization          (ICRA: dynamics randomization for sim-to-real)
%  - Longhini2025ReviewClothManipulation    (Annual Review of robotic cloth manipulation)
%  - Nealen2006DeformableModels             (STAR: deformable model families and tradeoffs)
%  - BaraffWitkin1998LargeStepsCloth        (SIGGRAPH: stable cloth simulation with larger steps)
%  - Bridson2002ClothCollisions             (SIGGRAPH: robust cloth collision/contact/friction)
%  - SifakisBarbic2012FEM                   (SIGGRAPH course: FEM simulation fundamentals)
%  - Jiang2016MPM                           (SIGGRAPH/TOG: MPM for continuum materials)
%  - muller2007pbd                          (JVCIR: Position Based Dynamics)
%  - macklin2016xpbd                        (MIG: XPBD compliance and Lagrange multipliers)
%  - FratarcangeliPellacini2013GPUPBD       (JGT: GPU PBD with graph coloring)
%  - Macklin2019SmallSteps                  (Small Steps in Physics Simulation: substepping insight)
%  - Chen2024VBD                            (TOG: Vertex Block Descent)
%  - newton_physics_doc                     (Newton Physics documentation)
%  - nvidia_newton_announcement_2025        (NVIDIA blog: announcing Newton)
%  - nvidia_isaaclab_newton_blog_2025       (NVIDIA blog: Isaac Lab + Newton cloth manipulation)
% ---------------------------------------------------------------------------

% ---------------------------------------------------------------------------
% BibTeX entries — add to bibliography/literature.bib:
%
% @online{omni_physx_overview,
%   author  = {{NVIDIA}},
%   title   = {Omni PhysX --- Overview},
%   year    = {2026},
%   url     = {https://docs.omniverse.nvidia.com/kit/docs/omni_physics/latest/extensions/runtime/source/omni.physx/docs/index.html},
%   urldate = {2026-03-03},
% }
%
% @online{openusd_usdphysics_schema,
%   author  = {{Pixar Animation Studios}},
%   title   = {UsdPhysics: {USD} Physics Schema},
%   year    = {2021},
%   url     = {https://openusd.org/docs/api/usd_physics_page_front.html},
%   urldate = {2026-03-03},
% }
%
% @online{isaac_sim_simulation_fundamentals_510,
%   author  = {{NVIDIA}},
%   title   = {Physics Simulation Fundamentals --- Isaac Sim 5.1.0 Documentation},
%   year    = {2026},
%   url     = {https://docs.isaacsim.omniverse.nvidia.com/5.1.0/physics/simulation_fundamentals.html},
%   urldate = {2026-03-03},
% }
%
% @online{isaac_sim_physics_resources_510,
%   author  = {{NVIDIA}},
%   title   = {Omniverse Physics and PhysX SDK Limitations --- Isaac Sim 5.1.0 Documentation},
%   year    = {2026},
%   url     = {https://docs.isaacsim.omniverse.nvidia.com/5.1.0/physics/physics_resources.html},
%   urldate = {2026-03-03},
% }
%
% @online{physx_welcome_513,
%   author  = {{NVIDIA}},
%   title   = {Welcome to PhysX --- PhysX 5.1.3 Documentation},
%   year    = {2023},
%   url     = {https://nvidia-omniverse.github.io/PhysX/physx/5.1.3/index.html},
%   urldate = {2026-03-03},
% }
%
% @online{physx_rigid_body_dynamics_511,
%   author  = {{NVIDIA}},
%   title   = {Rigid Body Dynamics --- PhysX 5.1.1 Documentation},
%   year    = {2022},
%   url     = {https://nvidia-omniverse.github.io/PhysX/physx/5.1.1/docs/RigidBodyDynamics.html},
%   urldate = {2026-03-03},
% }
%
% @online{physx_rigid_body_collision_510,
%   author  = {{NVIDIA}},
%   title   = {Rigid Body Collision --- PhysX 5.1.0 Documentation},
%   year    = {2022},
%   url     = {https://nvidia-omniverse.github.io/PhysX/physx/5.1.0/docs/RigidBodyCollision.html},
%   urldate = {2026-03-03},
% }
%
% @online{physx_material_api_510,
%   author  = {{NVIDIA}},
%   title   = {{PxMaterial} --- PhysX 5.1.0 API Reference},
%   year    = {2022},
%   url     = {https://nvidia-omniverse.github.io/PhysX/physx/5.1.0/_build/physx/latest/class_px_material.html},
%   urldate = {2026-03-03},
% }
%
% @online{physx_simulation_541,
%   author  = {{NVIDIA}},
%   title   = {Simulation --- PhysX 5.4.1 Documentation},
%   year    = {2024},
%   url     = {https://nvidia-omniverse.github.io/PhysX/physx/5.4.1/docs/Simulation.html},
%   urldate = {2026-03-03},
% }
%
% @online{physx_best_practices_541,
%   author  = {{NVIDIA}},
%   title   = {Best Practices Guide --- PhysX 5.4.1 Documentation},
%   year    = {2024},
%   url     = {https://nvidia-omniverse.github.io/PhysX/physx/5.4.1/docs/BestPractices.html},
%   urldate = {2026-03-03},
% }
%
% @online{physx_advanced_collision_detection_512,
%   author  = {{NVIDIA}},
%   title   = {Advanced Collision Detection --- PhysX 5.1.2 Documentation},
%   year    = {2022},
%   url     = {https://nvidia-omniverse.github.io/PhysX/physx/5.1.2/docs/AdvancedCollisionDetection.html},
%   urldate = {2026-03-03},
% }
%
% @online{physx_gpu_rigid_bodies_510,
%   author  = {{NVIDIA}},
%   title   = {Using GPU Rigid Bodies --- PhysX 5.1.0 Documentation},
%   year    = {2022},
%   url     = {https://nvidia-omniverse.github.io/PhysX/physx/5.1.0/docs/GPURigidBodies.html},
%   urldate = {2026-03-03},
% }
%
% @online{physx_geometry_541,
%   author  = {{NVIDIA}},
%   title   = {Geometry --- PhysX 5.4.1 Documentation},
%   year    = {2024},
%   url     = {https://nvidia-omniverse.github.io/PhysX/physx/5.4.1/docs/Geometry.html},
%   urldate = {2026-03-03},
% }
%
% @online{physx_geometry_530,
%   author  = {{NVIDIA}},
%   title   = {Geometry --- PhysX 5.3.0 Documentation},
%   year    = {2023},
%   url     = {https://nvidia-omniverse.github.io/PhysX/physx/5.3.0/docs/Geometry.html},
%   urldate = {2026-03-03},
% }
%
% @online{physx_articulations_513,
%   author  = {{NVIDIA}},
%   title   = {Articulations --- PhysX 5.1.3 Documentation},
%   year    = {2023},
%   url     = {https://nvidia-omniverse.github.io/PhysX/physx/5.1.3/docs/Articulations.html},
%   urldate = {2026-03-03},
% }
%
% @online{omni_physics_tendon_authoring,
%   author  = {{NVIDIA}},
%   title   = {Tendon Visual Authoring --- Omni Physics Documentation},
%   year    = {2025},
%   url     = {https://docs.omniverse.nvidia.com/kit/docs/omni_physics/latest/extensions/ux/source/omni.usdphysics.ui/docs/dev_guide/tendon_authoring.html},
%   urldate = {2026-03-03},
% }
%
% @online{isaaclab_sim_doc,
%   author  = {{NVIDIA}},
%   title   = {isaaclab.sim --- IsaacLab Documentation (Simulation Configuration and GPU Physics)},
%   year    = {2025},
%   url     = {https://isaac-sim.github.io/IsaacLab/main/source/api/lab/isaaclab.sim.html},
%   urldate = {2026-03-03},
% }
%
% @online{physx_particle_system_513,
%   author  = {{NVIDIA}},
%   title   = {Particle System --- PhysX 5.1.3 Documentation},
%   year    = {2023},
%   url     = {https://nvidia-omniverse.github.io/PhysX/physx/5.1.3/docs/ParticleSystem.html},
%   urldate = {2026-03-03},
% }
%
% @online{physx_pbd_particlesystem_api,
%   author  = {{NVIDIA}},
%   title   = {{PxPBDParticleSystem} --- PhysX API Reference},
%   year    = {2023},
%   url     = {https://nvidia-omniverse.github.io/PhysX/physx/5.3.1/_api_build/class_px_p_b_d_particle_system.html},
%   urldate = {2026-03-03},
% }
%
% @online{omni_physics_particles,
%   author  = {{NVIDIA}},
%   title   = {Particles --- Omni Physics Documentation},
%   year    = {2025},
%   url     = {https://docs.omniverse.nvidia.com/kit/docs/omni_physics/latest/dev_guide/particles/particles.html},
%   urldate = {2026-03-03},
% }
%
% @inproceedings{Makoviychuk2021IsaacGym,
%   author    = {Makoviychuk, Viktor and Wawrzyniak, Lukasz and Guo, Yunrong and Lu, Michelle and Storey, Kier and Macklin, Miles and Hoeller, David and Rudin, Nikita and Allshire, Arthur and Handa, Ankur and State, Gavriel},
%   title     = {Isaac Gym: High Performance GPU-Based Physics Simulation for Robot Learning},
%   booktitle = {Proceedings of the Neural Information Processing Systems Track on Datasets and Benchmarks},
%   year      = {2021},
%   url       = {https://datasets-benchmarks-proceedings.neurips.cc/paper/2021/file/28dd2c7955ce926456240b2ff0100bde-Paper-round2.pdf},
% }
%
% @inproceedings{Jakobi1995RealityGap,
%   author    = {Jakobi, Nick and Husbands, Phil and Harvey, Inman},
%   title     = {Noise and the Reality Gap: The Use of Simulation in Evolutionary Robotics},
%   booktitle = {Advances in Artificial Life (ECAL 1995)},
%   series    = {Lecture Notes in Computer Science},
%   volume    = {929},
%   pages     = {704--720},
%   year      = {1995},
%   publisher = {Springer},
%   doi       = {10.1007/3-540-59496-5_337},
% }
%
% @article{Salvato2021RealityGapSurvey,
%   author  = {Salvato, Erica and Fenu, Gianfranco and Medvet, Eric and Pellegrino, Felice Andrea},
%   title   = {Crossing the Reality Gap: A Survey on Sim-to-Real Transferability of Robot Controllers in Reinforcement Learning},
%   journal = {IEEE Access},
%   volume  = {9},
%   pages   = {153171--153187},
%   year    = {2021},
%   doi     = {10.1109/ACCESS.2021.3126658},
% }
%
% @inproceedings{Peng2018DynamicsRandomization,
%   author    = {Peng, Xue Bin and Andrychowicz, Marcin and Zaremba, Wojciech and Abbeel, Pieter},
%   title     = {Sim-to-Real Transfer of Robotic Control with Dynamics Randomization},
%   booktitle = {2018 IEEE International Conference on Robotics and Automation (ICRA)},
%   year      = {2018},
%   pages     = {3803--3810},
%   doi       = {10.1109/ICRA.2018.8460528},
% }
%
% @article{Longhini2025ReviewClothManipulation,
%   author  = {Longhini, Alberta and Wang, Yufei and Garcia-Camacho, Irene and Blanco-Mulero, David and Moletta, Marco and Welle, Michael and Aleny{\`a}, Guillem and Yin, Hang and Erickson, Zackory and Held, David and Borr{\`a}s, J{\'u}lia and Kragic, Danica},
%   title   = {A Review of Robotic Cloth Manipulation},
%   journal = {Annual Review of Control, Robotics, and Autonomous Systems},
%   year    = {2025},
%   doi     = {10.1146/annurev-control-022723-033252},
% }
%
% @article{Nealen2006DeformableModels,
%   author  = {Nealen, Andrew and M{\"u}ller, Matthias and Keiser, Richard and Boxerman, Eddy and Carlson, Mark},
%   title   = {Physically Based Deformable Models in Computer Graphics},
%   journal = {Computer Graphics Forum},
%   volume  = {25},
%   number  = {4},
%   pages   = {809--836},
%   year    = {2006},
%   doi     = {10.1111/j.1467-8659.2006.01000.x},
% }
%
% @inproceedings{BaraffWitkin1998LargeStepsCloth,
%   author    = {Baraff, David and Witkin, Andrew},
%   title     = {Large Steps in Cloth Simulation},
%   booktitle = {Proceedings of the 25th Annual Conference on Computer Graphics and Interactive Techniques (SIGGRAPH '98)},
%   year      = {1998},
%   pages     = {43--54},
%   doi       = {10.1145/280814.280821},
% }
%
% @inproceedings{Bridson2002ClothCollisions,
%   author    = {Bridson, Robert and Fedkiw, Ronald and Anderson, John},
%   title     = {Robust Treatment of Collisions, Contact and Friction for Cloth Animation},
%   booktitle = {Proceedings of the 29th Annual Conference on Computer Graphics and Interactive Techniques (SIGGRAPH '02)},
%   year      = {2002},
%   pages     = {594--603},
%   doi       = {10.1145/566570.566623},
% }
%
% @inproceedings{SifakisBarbic2012FEM,
%   author    = {Sifakis, Eftychios and Barbi{\v{c}}, Jernej},
%   title     = {{FEM} Simulation of 3D Deformable Solids: A Practitioner's Guide to Theory, Discretization and Model Reduction},
%   booktitle = {ACM SIGGRAPH 2012 Courses},
%   year      = {2012},
%   doi       = {10.1145/2343483.2343501},
% }
%
% @article{Jiang2016MPM,
%   author  = {Jiang, Chenfanfu and Schroeder, Craig and Teran, Joseph and Stomakhin, Alexey and Selle, Andrew},
%   title   = {The Material Point Method for Simulating Continuum Materials},
%   journal = {ACM Transactions on Graphics},
%   volume  = {35},
%   number  = {4},
%   year    = {2016},
%   doi     = {10.1145/2897826.2927348},
%   url     = {https://alexey.stomakhin.com/research/siggraph2016_mpm.pdf},
% }
%
% @article{muller2007pbd,
%   author  = {M{\"u}ller, Matthias and Heidelberger, Bruno and Hennix, Marcus and Ratcliff, John},
%   title   = {Position Based Dynamics},
%   journal = {Journal of Visual Communication and Image Representation},
%   volume  = {18},
%   number  = {2},
%   pages   = {109--118},
%   year    = {2007},
%   doi     = {10.1016/j.jvcir.2007.01.005},
% }
%
% @inproceedings{macklin2016xpbd,
%   author    = {Macklin, Miles and M{\"u}ller, Matthias and Chentanez, Nuttapong},
%   title     = {XPBD: Position-Based Simulation of Compliant Constrained Dynamics},
%   booktitle = {Proceedings of the 9th International Conference on Motion in Games (MIG '16)},
%   year      = {2016},
%   pages     = {49--54},
%   doi       = {10.1145/2994258.2994272},
% }
%
% @article{FratarcangeliPellacini2013GPUPBD,
%   author  = {Fratarcangeli, Marco and Pellacini, Fabio},
%   title   = {A GPU-Based Implementation of Position Based Dynamics for Interactive Deformable Bodies},
%   journal = {Journal of Graphics Tools},
%   volume  = {17},
%   number  = {3},
%   pages   = {59--66},
%   year    = {2013},
%   doi     = {10.1080/2165347X.2015.1030525},
% }
%
% @online{Macklin2019SmallSteps,
%   author  = {Macklin, Miles and Storey, Kier and Lu, Michelle and Terdiman, Pierre and Chentanez, Nuttapong and Jeschke, Stefan and M{\"u}ller, Matthias},
%   title   = {Small Steps in Physics Simulation},
%   year    = {2019},
%   url     = {https://mmacklin.com/smallsteps.pdf},
%   urldate = {2026-03-03},
% }
%
% @article{Chen2024VBD,
%   author  = {Chen, Anka He and Liu, Ziheng and Yang, Yin and Yuksel, Cem},
%   title   = {Vertex Block Descent},
%   journal = {ACM Transactions on Graphics},
%   volume  = {43},
%   number  = {4},
%   pages   = {1--16},
%   year    = {2024},
%   doi     = {10.1145/3658179},
%   url     = {https://arxiv.org/abs/2403.06321},
% }
%
% @online{newton_physics_doc,
%   author  = {{Newton Physics Contributors}},
%   title   = {Newton Physics Documentation},
%   year    = {2025},
%   url     = {https://newton-physics.github.io/newton/},
%   urldate = {2026-03-03},
% }
%
% @online{nvidia_newton_announcement_2025,
%   author  = {{NVIDIA}},
%   title   = {Announcing Newton, an Open-Source Physics Engine for Robotics Simulation},
%   year    = {2025},
%   url     = {https://developer.nvidia.com/blog/announcing-newton-an-open-source-physics-engine-for-robotics-simulation/},
%   urldate = {2026-03-03},
% }
%
% @online{nvidia_isaaclab_newton_blog_2025,
%   author  = {{NVIDIA}},
%   title   = {Train a Quadruped Locomotion Policy and Simulate Cloth Manipulation with NVIDIA Isaac Lab and Newton},
%   year    = {2025},
%   url     = {https://developer.nvidia.com/blog/train-a-quadruped-locomotion-policy-and-simulate-cloth-manipulation-with-nvidia-isaac-lab-and-newton/},
%   urldate = {2026-03-03},
% }
% ---------------------------------------------------------------------------